\Needspace{11\baselineskip}
\SourceLeaf{273}
\begin{center}
\textit{Composition of genera.}
\end{center}

\Article{246}
\textsc{Problem.} \textit{Given any two primitive forms $f,f'$, whose composition gives $F$, determine from the genera to which $f,f'$ belong the genus to which $F$ belongs.}

\textit{Solution.} I. First consider the case where at least one of the forms $f,f'$, say the first, is properly primitive, and denote the determinants of the forms $f,f',F$ by $d,d',D$. Then $D$ is the greatest common divisor of $dm'm'$, $d'$, where $m'$ is either $1$ or $2$, according as the form $f'$ is properly or improperly primitive; $F$ belongs to the properly primitive order in the former case and to the improperly primitive order in the latter. Now the genus of the form $F$ is determined by its individual characters, with respect both to each odd prime divisor of $D$ and, in certain cases, to the numbers $4$ or $8$. We must therefore determine each of these.

$1^{\circ}$. If $p$ is any odd prime divisor of $D$, it must also divide $d,d'$, so the relations of the forms $f,f'$ to $p$ also occur among their characters. Now if $f$ can represent the number $a$, and $f'$ the number $a'$, the product $aa'$ can be represented by $F$. So if both $f$ and $f'$ can represent quadratic residues modulo $p$ (not divisible by $p$), $F$ can also represent quadratic residues modulo $p$; \textit{that is,}\ if both $f,f'$ have the character $Rp$, the form $F$ has the same character. Likewise $F$ has the character $Rp$ if both $f,f'$ have the character $Np$; by contrast, $F$ has the character $Np$ if one of the forms $f,f'$ has $Rp$ and the other $Np$.

$2^{\circ}$. If the complete character of the form $F$ includes a relation to the number $4$, such a relation must also occur among the characters of the forms $f$, $f'$. For this happens only when $D$ is $\equiv0$ or $\equiv3$ (modulo $4$). When $D$ is divisible by $4$, both $dm'm'$ and $d'$ are also divisible by $4$. It follows at once that $f'$ cannot be improperly primitive, and therefore $m'=1$; thus both $d$ and $d'$ are divisible by $4$, and the character of each includes a relation to $4$. When $D\equiv3$ (modulo $4$), $D$ divides $d,d'$; the quotients are squares, so $d,d'$ must also be either $\equiv0$ or $\equiv3$ (modulo $4$), and among the \SourceLeaf{274}characters of $f,f'$ a relation to $4$ occurs. It follows, just as in ($1^{\circ}$), that the form $F$ has the character $1,4$ if either both $f,f'$ have $1,4$ or both have $3,4$; by contrast, the form $F$ has the character $3,4$ if one of the forms $f,f'$ has $1,4$ and the other $3,4$.

$3^{\circ}$. When $D$ is divisible by $8$, $d'$ is too; so $f'$ is certainly properly primitive, $m'=1$, and $d$ is also divisible by $8$. Thus among the characters of the form $F$, one of the characters $1,8$; $3,8$; $5,8$; $7,8$ can occur only if the characters of both forms $f$, $f'$ also include such a relation to $8$. As before, we easily confirm that the form $F$ has the character $1,8$ if $f$ and $f'$ have the same character with respect to $8$; the form $F$ has the character $3,8$ if one of the forms $f,f'$ has $1,8$ and the other $3,8$, or one has $5,8$ and the other $7,8$; $F$ has $5,8$ if $f,f'$ have $1,8$ and $5,8$, or $3,8$ and $7,8$; $F$ has $.7,8$ if $f$ and $f'$ have either $1,8$ and $7,8$, or $3,8$ and $5,8$.

$4^{\circ}$. When $D\equiv2$ (modulo $8$), $d'$ is either $\equiv0$ or $\equiv2$ (modulo $8$), so $m'=1$, and therefore $d$ is also either $\equiv0$ or $\equiv2$ (modulo $8$). But $d,d'$ cannot both be divisible by $8$, since $D$ is their \textit{greatest} common divisor. Thus one of the characters $1\text{ \textit{and} }7,8$; $3\text{ \textit{and} }5,8$ can be assigned to the form $F$ only when either both forms $f,f'$ have one of those characters, or one has one of them and the other one of the characters $1,8$; $3,8$; $5,8$; $7,8$. It is easy to deduce that the character of the form $F$ is determined by the following table, where the character in the margin belongs to one of the forms $f,f'$ and the character at the top to the other:
\[
\begin{array}{||c||c|c||}
\hline\hline
 & 1\text{ \textit{and} }7,8 & 3\text{ \textit{and} }5,8 \\
 & \text{or }1,8 & \text{or }3,8 \\
 & \text{or }7,8 & \text{or }5,8 \\
\hline\hline
1\text{ \textit{and} }7,8 & 1\text{ \textit{and} }7,8 & 3\text{ \textit{and} }5,8 \\
\hline
3\text{ \textit{and} }5,8 & 3\text{ \textit{and} }5,8 & 1\text{ \textit{and} }7.8 \\
\hline\hline
\end{array}
\]

$5^{\circ}$. In the same way we prove that the form $F$ cannot be assigned either of the characters $1\text{ \textit{and} }3,8$; $5\text{ \textit{and} }7,8$ unless at least one of the forms $f,f'$ also has one of those characters and the other either one of them or one of the characters $1,8$; $3,8$; $5,8$; $7,8$.
\SourceLeaf{275}
The character of the form $F$ is determined by the following table, whose margin and top give the characters of the forms $f,f'$:
\[
\begin{array}{||c||c|c||}
\hline\hline
 & 1\text{ \textit{and} }3,8 & 5\text{ \textit{and} }7,8 \\
 & \text{or }1,8 & \text{or }5,8 \\
 & \text{or }3,8 & \text{or }7,8 \\
\hline\hline
1\text{ \textit{and} }3,8 & 1\text{ \textit{and} }3,8 & 5\text{ \textit{and} }7,8 \\
\hline
5\text{ \textit{and} }7,8 & 5\text{ \textit{and} }7,8 & 1\text{ \textit{and} }3.8 \\
\hline\hline
\end{array}
\]

II. If both forms $f,f'$ are improperly primitive, $D$ is the greatest common divisor of $4d.\ 4d'$, or $\frac{1}{4}D$ is the greatest common divisor of $d.\ d'$. It follows easily that $d$, $d'$ and $\frac{1}{4}D$ are all $\equiv1$ (modulo $4$). Setting $F=(A,B,C)$, the greatest common divisor of $A,B,C$ is $=2$. and the greatest common divisor of $A,2B,C$ is $4$. Thus $F$ is a form derived from the improperly primitive form $(\frac{1}{2}A,\frac{1}{2}B,\frac{1}{2}C)$, whose determinant is $\frac{1}{4}D$ and whose genus determines the genus of the form $F$. As that form is improperly primitive, its character includes no relations to $4$ or $8$, only relations to the individual odd prime divisors of $\frac{1}{4}D$. Now all these divisors clearly also divide $d,d'$, and half of any product of two factors. one representable by $f$ and the other by $f'$, can be represented by the form $(\frac{1}{2}A.\ \frac{1}{2}B,\frac{1}{2}C)$. It is therefore easy to see that the character of this form with respect to any odd prime $p$ dividing $\frac{1}{4}D$ is $Rp$, \textit{both} if $2Rp$ and the forms $f,f'$ have the same character with respect to $p$. \textit{and} if $2Np$ and the characters of the forms $f,f'$ with respect to $p$ are opposite; by contrast, the character of that form is $Np$, \textit{both} if $f,f'$ have the same character with respect to $p$ and $2Np$, \textit{and} if $f,f'$ have opposite characters and $2Rp$.

\Needspace{8\baselineskip}
\Article{247}
The solution of the preceding problem shows that if $g$ is a primitive form of the same order and genus as $f$, and $g'$ a primitive form of the same order and genus as $f'$, the form composed from $g$ and $g'$ belongs to the same genus as the form composed from $f$ and $f'$. This gives the meaning of a \textit{genus}
\SourceLeaf{276}
\textit{composed} from two other genera (or even more). Further, the same solution shows that if $f,f'$ have the same determinant, $f$ is a form of the principal genus and $F$ is composed from $f$ and $f'$, then $F$ belongs to the same genus as $f'$. So the principal genus can always be left out when composing it with other genera of the same determinant. But if, with the other conditions unchanged, $f$ does not belong to the principal genus and $f'$ is a primitive form, $F$ certainly belongs to a different genus from $f'$. Finally, if $f,f'$ are properly primitive forms of the same genus, $F$ belongs to the principal genus; but if $f,f'$ are both properly primitive with the same determinant but of different genera, $F$ cannot belong to the principal genus. Thus, when any properly primitive form is composed \textit{with itself}, the resulting form, which is also properly primitive with the same determinant, must belong to the principal genus.

\Needspace{8\baselineskip}
\Article{248}
\textsc{Problem.} \textit{Given any two forms $f,f'$, from which $F$ is composed, determine from the genera of the forms $f,f'$ the genus of the form $F$.}

\textit{Solution.} Let $f=(a,b,c)$, $f'=(a',b',c')$, $F=(A,B,C)$, and let $\mathfrak{m}$ be the greatest common divisor of $a,b,c$, and $\mathfrak{m}'$ the greatest common divisor of $a',b',c'$, so that $f,f'$ are derived from the primitive forms $(\frac{a}{\mathfrak{m}},\allowbreak \frac{b}{\mathfrak{m}},\allowbreak \frac{c}{\mathfrak{m}})$, $(\frac{a'}{\mathfrak{m}'},\allowbreak \frac{b'}{\mathfrak{m}'},\allowbreak \frac{c'}{\mathfrak{m}'})$, which we denote by $\mathfrak{f}$, $\mathfrak{f}'$ respectively. Now, if at least one of the forms $\mathfrak{f}$, $\mathfrak{f}'$ is properly primitive, the greatest common divisor of $A,B,C$ is $\mathfrak{m}\mathfrak{m}'$, so $F$ is derived from the primitive form $(\frac{A}{\mathfrak{m}\mathfrak{m}'},\allowbreak \frac{B}{\mathfrak{m}\mathfrak{m}'},\allowbreak \frac{C}{\mathfrak{m}\mathfrak{m}'})\ldots\mathfrak{F}$; thus the genus of the form $F$ depends on the genus of the form $\mathfrak{F}$. But it is easy to see that $\mathfrak{F}$ is transformed into $\mathfrak{f}\mathfrak{f}'$ by the same substitution that transforms $F$ into $ff'$, so $\mathfrak{F}$ is composed from $\mathfrak{f}$, $\mathfrak{f}'$, and its genus can be determined by the problem of Article~246. --- But if both $\mathfrak{f}$, $\mathfrak{f}'$ are improperly primitive, the greatest common divisor of $A,B,C$ is $2\mathfrak{m}\mathfrak{m}'$, and the form $\mathfrak{F}$ is still composed from $\mathfrak{f}$, $\mathfrak{f}'$ and is clearly derived from the properly primitive form $(\frac{A}{2\mathfrak{m}\mathfrak{m}'},\allowbreak \frac{B}{2\mathfrak{m}\mathfrak{m}'},\allowbreak \frac{C}{2\mathfrak{m}\mathfrak{m}'})$. The genus of this form can therefore be determined by Article~246; and since $F$ is derived from the same form, its genus follows at once.

This solution shows that the theorem set out in the preceding article for primitive forms, namely \textit{if $f',g'$ belong to the same genera as $f,g$ respectively, the form composed from $f',g'$ belongs to the same genus as the form composed from $f,g$}, holds in general for any forms.
